\chapter*{Preface}

This book is a hands-on guide to building a large language model from the ground up.
It covers every major component---from tokenizing raw text, through attention and the
transformer architecture, to pretraining, loading pretrained weights, fine-tuning,
and modern architectural innovations---with enough depth that you can implement each
piece yourself.

\section*{Who This Book Is For}

You should be comfortable writing Python and have some experience with PyTorch,
including tensors, automatic differentiation, and \texttt{nn.Module}.
We assume familiarity with basic linear algebra (matrix multiplication, dot products)
and calculus (gradients, the chain rule). No prior experience with language models
or transformers is required; we build every concept from scratch.

\section*{Our Reference Model: GPT-2}

Throughout this book, we focus on implementing a GPT-2--class language model.
Released by OpenAI in 2019, GPT-2 is an ideal learning target for several
reasons:

\begin{itemize}
    \item \textbf{Simplicity}: GPT-2 uses a clean, decoder-only transformer
          architecture without the encoder--decoder complexity of the original
          Transformer or the architectural additions of later models. Every
          component---embeddings, attention, layer normalization, feed-forward
          layers---can be understood and implemented in isolation.
    \item \textbf{Pretrained weights}: OpenAI published trained GPT-2 models on
          HuggingFace, so we can load real weights (Chapters~5--6) and verify
          our implementation produces coherent text, without needing to train
          from scratch on a GPU cluster.
    \item \textbf{Scale}: At 124M parameters (GPT-2 Small), the model is small
          enough to train and experiment with on a consumer GPU, yet large
          enough to demonstrate the key engineering challenges of real LLMs.
    \item \textbf{Foundation for modern models}: GPT-3, LLaMA, Mistral, and
          most contemporary LLMs are direct descendants of the GPT-2
          architecture. The modifications they introduce---RMSNorm, rotary
          positional embeddings, grouped query attention, SwiGLU---are covered
          in Chapter~7, building directly on the GPT-2 baseline.
\end{itemize}

When we refer to ``the model'' or ``our implementation'' without
qualification, we mean GPT-2 Small (12 layers, 768-dimensional embeddings,
12 attention heads). Larger GPT-2 variants (Medium, Large, XL) share the
same architecture with different hyperparameters.

\section*{Pseudocode Convention}

Throughout this book, algorithms are presented in Python-like pseudocode with type
hints and tensor shape annotations. For example:

\begin{pseudocode}[Example: Shape Annotations]
def some_function(x: Tensor) -> Tensor:  # x: [batch, seq_len, embed_dim]
    y = linear(x)  # [batch, seq_len, out_dim]
    return y        # [batch, seq_len, out_dim]
\end{pseudocode}

Shape annotations in comments track how tensor dimensions change through each
operation. Where abbreviations are necessary for layout reasons, a legend is
provided at the top of the code block explaining each abbreviation.

\section*{Chapter Overview}

\begin{enumerate}
    \item \textbf{Tokenization} --- Turning raw text into sequences of integers
    \item \textbf{Attention} --- The mechanism that lets tokens interact with each other
    \item \textbf{The GPT Architecture} --- Assembling the full transformer decoder
    \item \textbf{Pretraining} --- Training the model with next-token prediction
    \item \textbf{Loading Pretrained Weights} --- Reusing weights from HuggingFace
    \item \textbf{Fine-Tuning} --- Adapting a pretrained model for classification and instruction following
    \item \textbf{Modern Architecture Deep Dives} --- RMSNorm, RoPE, SwiGLU, and GQA
\end{enumerate}

Each chapter ends with review questions to test your understanding and hands-on
coding problems to strengthen your implementation skills.